% 数学ノート集・共通テンプレート v2（積分論2026年度第1部の版面を参考）
% A5判。章内で定義・定理・例・注意等の番号を共通化。
% コンパイル： platex -interaction=nonstopmode math_note_template_v2.tex (2回)
%              dvipdfmx math_note_template_v2.dvi
\documentclass[a5paper,10pt,onecolumn]{jsarticle}
\usepackage[dvipdfmx,top=22mm,bottom=20mm,inner=19mm,outer=18mm,headheight=13pt,headsep=6mm]{geometry}
\usepackage{amsmath,amssymb,amsthm,mathtools}
\usepackage{enumitem,fancyhdr,titlesec}
\usepackage[dvipdfmx,bookmarks=true,colorlinks=false]{hyperref}
\usepackage{pxjahyper}

% ----- 差し替える文書情報 -----
\newcommand{\NoteTitle}{Godelの不完全性定理}
\newcommand{\NoteSeries}{数学ノート集}

% ----- 文字・組版 -----
% 参考PDFの質素な明朝＋Computer Modernの印象を引き継ぎながら，
% A5判で窮屈にならないよう本文の行間を少し確保する．
\linespread{1.13}
\setlength{\parindent}{1zw}
\setlength{\parskip}{0pt}
\setlength{\abovedisplayskip}{6pt plus 2pt minus 2pt}
\setlength{\belowdisplayskip}{6pt plus 2pt minus 2pt}
\setlength{\abovedisplayshortskip}{3pt plus 2pt minus 1pt}
\setlength{\belowdisplayshortskip}{4pt plus 2pt minus 1pt}
\setlist[enumerate]{topsep=3pt,itemsep=1pt,parsep=0pt,leftmargin=2.4zw}
\titleformat{\section}{\normalfont\normalsize\bfseries}{\thesection}{0.8zw}{}
\titlespacing*{\section}{0pt}{2.0ex plus .4ex}{.9ex}
\titleformat{\subsection}{\normalfont\normalsize\bfseries}{\thesubsection}{0.7zw}{}
\titlespacing*{\subsection}{0pt}{1.5ex plus .3ex}{.6ex}
\raggedbottom

% ----- 柱とページ番号：ページ上部、罫線なし -----
\fancypagestyle{mathnote}{%
  \fancyhf{}%
  \fancyhead[L]{\footnotesize \NoteSeries\ （\NoteTitle）}%
  \fancyhead[R]{\footnotesize \thepage}%
  \renewcommand{\headrulewidth}{0pt}%
  \renewcommand{\footrulewidth}{0pt}%
}
\pagestyle{mathnote}

% ----- 節内通し番号の定義・定理群 -----
\newtheoremstyle{mathnotetheorem}
 {4.5pt}{4.5pt}{\normalfont}{}%
 {\bfseries}{}{0.65zw}%
 {\thmname{#1}\thmnumber{ #2}\thmnote{（#3）}}
\theoremstyle{mathnotetheorem}
\newtheorem{theorem}{定理}[section]
\newtheorem{definition}[theorem]{定義}
\newtheorem{lemma}[theorem]{補題}
\newtheorem{proposition}[theorem]{命題}
\newtheorem{corollary}[theorem]{系}
\newtheorem{example}[theorem]{例}
\newtheorem{remark}[theorem]{注意}
\renewcommand{\proofname}{証明}
\renewcommand{\qedsymbol}{$\square$}
\makeatletter
\renewenvironment{proof}[1][\proofname]{%
  \par\pushQED{\qed}\normalfont\topsep4\p@\@plus2\p@\relax
  \trivlist\item[\hskip\labelsep\bfseries #1]\ignorespaces
}{\popQED\endtrivlist\@endpefalse}
\makeatother

% ----- 使うノートで適宜拡張する記号 -----
\newcommand{\R}{\mathbb{R}}
\newcommand{\norm}[1]{\lVert #1\rVert}
\newcommand{\abs}[1]{\lvert #1\rvert}

\begin{document}
\section*{読了後にできること}
本稿を読むための前提は，証明可能性述語と対角化補題．
本稿の具体的な到達課題は，第一不完全性定理の不完全命題を構成する論理的矛盾を再現する．

\section{形式的算術と計算可能な理論}

本稿では古典的一階論理の健全性・完全性定理，
原始再帰関数，Peano算術 $\mathrm{PA}$ の公理を前提とする．
全ての構文対象は有限文字列として符号化する．
理論の無矛盾性と標準モデルでの真理性を区別する．

\begin{definition}
一階理論 $T$ が\emph{計算可能に公理化可能}であるとは，
その公理集合が再帰的可算であることをいう．
$T$ が無矛盾であるとは，ある文 $\varphi$ について
$T\vdash\varphi$ と $T\vdash\neg\varphi$ が
同時に成り立たないことをいう．
\end{definition}

\begin{definition}
$T$ が\emph{$\omega$-無矛盾}であるとは，
一変数論理式 $\varphi(x)$ に対して
$T\vdash\exists x\,\varphi(x)$ が成り立つ一方，
全ての自然数 $n$ に対して
$T\vdash\neg\varphi(\overline n)$ が成り立つ，
という状況が存在しないことをいう．
ここで $\overline n$ は自然数 $n$ の数詞である．
\end{definition}

\begin{proposition}
$\omega$-無矛盾な理論は無矛盾である．
\end{proposition}
\begin{proof}
もし $T$ が矛盾すれば任意の文が証明可能となる．
従って $\exists x\,(x=x)$ も，
全ての $n$ に対する $\neg(\overline n=\overline n)$ も
証明可能であり，$\omega$-無矛盾性に反する．
\end{proof}

\section{自然数による構文の符号化}

\begin{definition}
有限文字列の各文字へ正整数のコードを割り当てる．
文字コード $a_0,\dots,a_{r-1}$ からなる文字列には，
$r$ の記録と
\[
 \prod_{i=0}^{r-1}p_{i+1}^{a_i+1}
\]
を用いた一意な自然数コードを割り当てる．
ここで $p_i$ は $i$ 番目の素数である．
このようなコードをGödel番号という．
\end{definition}

\begin{lemma}[構文操作の計算可能性]\label{lem:syntax}
有限文字列の連結，長さの取得，式の正当性判定，
自由変数の検出，数詞の生成，
代入捕獲を避けた項の代入，証明列の一行の検査は
原始再帰的な符号操作として実現できる．
\end{lemma}
\begin{proof}
文字列の長さと各位置の文字は素因数の指数を
有界探索することで得られる．
$p_i$ の生成も，$i$ 個の素数が見つかるまでの
探索を原始再帰的な既知の上界の内側に制限すれば
原始再帰的に行える．
連結と部分列の取得は有限個の位置の文字コードを
組み替える操作なので原始再帰的である．
項や論理式の構文解析は，
入力文字列の長さ以下の有限個の構成木を調べ，
各節点の文法規則を判定する有界探索によって実現できる．
代入は構成木をたどり，束縛変数の名称を必要に応じて
新しい変数に変更してから有限文字列を置き換える操作である．
証明の一行が命題論理の公理例，量化公理例，
理論の公理，または先行行からの推論規則の結果であることも
有限個の行を調べることで判定できる．
理論の公理集合が再帰的可算の場合には，
各公理行へ列挙過程で現れた段階の証明書を付すことで，
証明書付き証明の検査を原始再帰的にできる．
従って述べた操作はいずれも原始再帰的に符号化できる．
\end{proof}

\begin{lemma}[Gödelの $\beta$-関数]\label{lem:beta}
任意の自然数の有限列 $a_0,\dots,a_r$ に対し，
自然数 $u,v$ が存在して
\[
 a_i=\operatorname{rem}(u,1+(i+1)v)
 \qquad(0\le i\le r)
\]
となる．
ここで $\operatorname{rem}$ は非負の除法の余りである．
\end{lemma}
\begin{proof}
$v$ を $(r+1)!$ の倍数で，
全ての $a_i$ より大きい正整数となるように取る．
$m_i=1+(i+1)v$ と置く．
$i<j$ のとき，$m_i,m_j$ の公約数 $d$ は
\[
 d\mid m_j-m_i=(j-i)v,\qquad
 \gcd(d,v)=1
\]
を満たすので $d\mid j-i$ となる．
一方 $v$ は $j-i$ の倍数だから
$m_i\equiv1\pmod{j-i}$ であり，
$d\mid m_i$ と $d\mid j-i$ から $d=1$ を得る．
よって $m_0,\dots,m_r$ は互いに素である．
中国剰余定理により
$u\equiv a_i\pmod{m_i}$ を同時に満たす $u$ が存在する．
$0\le a_i<m_i$ なのでその剰余はちょうど $a_i$ となる．
\end{proof}

\section{原始再帰関数の算術内での表現}

\begin{definition}
自然数上の関数 $f$ が $\mathrm{PA}$ で
\emph{強く表現可能}であるとは，
論理式 $F(\boldsymbol x,y)$ が存在して，
各自然数の組 $\boldsymbol n$ について
\[
 \mathrm{PA}\vdash F(\overline{\boldsymbol n},
                     \overline{f(\boldsymbol n)}),
\]
\[
 \mathrm{PA}\vdash
 \forall y\bigl(F(\overline{\boldsymbol n},y)
 \leftrightarrow y=\overline{f(\boldsymbol n)}\bigr)
\]
を満たすことをいう．
\end{definition}

\begin{theorem}[原始再帰関数の表現可能性]
\label{thm:representability}
全ての原始再帰関数は $\mathrm{PA}$ で強く表現可能である．
原始再帰的関係 $R$ に対しては，
真な数値例は $\mathrm{PA}$ で証明可能であり，
偽な数値例の否定も $\mathrm{PA}$ で証明可能であるような
論理式 $\rho_R$ を取れる．
\end{theorem}
\begin{proof}
原始再帰関数は，零関数，後続関数，射影関数から出発し，
合成と原始再帰の操作を有限回繰り返して得られる．
零・後続・射影のグラフは等号の式で表せるから，
各数値例の一意性も算術の公理から証明できる．

合成 $f(\boldsymbol x)
 =g(h_1(\boldsymbol x),\dots,h_m(\boldsymbol x))$
については，$g,h_i$ のグラフを表す論理式を用い
\[
 F(\boldsymbol x,y)\equiv
 \exists z_1\cdots z_m
 \left(\bigwedge_iH_i(\boldsymbol x,z_i)
 \land G(z_1,\dots,z_m,y)\right)
\]
と置く．
帰納法の仮定から各数値入力で $z_i$ と $y$ が
一意に指定されることを $\mathrm{PA}$ 内で証明できる．

原始再帰
\[
 f(0,\boldsymbol x)=g(\boldsymbol x),\qquad
 f(t+1,\boldsymbol x)
 =h(t,f(t,\boldsymbol x),\boldsymbol x)
\]
の場合には，
補題\ref{lem:beta}の余り表示を用いて
有限計算列
$a_0,\dots,a_t$ を二つの自然数 $u,v$ で符号化する．
$\beta(u,v,i)=\operatorname{rem}(u,1+(i+1)v)$ と書き，
\[
 \begin{aligned}
 F(t,\boldsymbol x,y)\equiv
 \exists u\,\exists v\ \bigl[&
 \beta(u,v,0)=g(\boldsymbol x)\ \land\\
 &\forall i<t\ 
 \beta(u,v,i+1)
 =h(i,\beta(u,v,i),\boldsymbol x)\ \land\\
 &\beta(u,v,t)=y\bigr]
 \end{aligned}
\]
を考える．
ここで $g,h$ の出現箇所は，帰納法で得た
それぞれのグラフ公式によって置き換える．
各標準数値の入力について，
計算列は有限個の具体的な値からなり，
補題\ref{lem:beta}の中国剰余構成は算術内で確認できる．
$\mathrm{PA}$ の帰納法公理によって，
この再帰関係を満たす計算列の各時刻の値は一意であると
証明できる．
従って $F$ は各標準数値入力で一意の出力を持つことが
証明可能である．
これで原始再帰関数の構成に関する帰納法が完了する．
関係 $R$ はその特性関数を用いればよい．
真偽の数値例について特性関数の値が0または1であることを
算術が証明できるため，後半が従う．
\end{proof}

\section{証明可能性述語}

\begin{definition}
計算可能に公理化可能な理論
$T\supseteq\mathrm{PA}$ を固定する．
証明書付き証明のコード $p$ が文のコード $x$ の
正しい $T$-証明であるという原始再帰的関係を
$\operatorname{Proof}_T(p,x)$ と書く．
定理\ref{thm:representability}により，
これを算術の論理式で表せる．
\[
 \operatorname{Prov}_T(x)
 \equiv\exists p\,\operatorname{Proof}_T(p,x)
\]
を証明可能性述語という．
\end{definition}

\begin{lemma}[数値例の検査]\label{lem:proof-numeral}
標準自然数 $p,x$ について，
$p$ が $x$ の正しい証明コードなら
\[
 \mathrm{PA}\vdash
 \operatorname{Proof}_T(\overline p,\overline x),
\]
そうでなければ
\[
 \mathrm{PA}\vdash
 \neg\operatorname{Proof}_T(\overline p,\overline x).
\]
特に $T\vdash\varphi$ なら
$\mathrm{PA}\vdash
 \operatorname{Prov}_T(\overline{\ulcorner\varphi\urcorner})$
となる．
\end{lemma}
\begin{proof}
$\operatorname{Proof}_T$ は原始再帰的関係なので，
定理\ref{thm:representability}から
全ての標準数値例の真偽が算術で証明可能である．
$T\vdash\varphi$ なら正しい証明コード $p$ が存在する．
その数値例の真を証明した後，存在量化を導入すると
$\operatorname{Prov}_T$ の証明を得る．
\end{proof}

\section{対角化補題}

\begin{lemma}[対角化補題]\label{lem:diagonal}
任意の一変数算術論理式 $\theta(x)$ に対し，
ある算術文 $G$ が存在して
\[
 \mathrm{PA}\vdash
 G\leftrightarrow
 \theta(\overline{\ulcorner G\urcorner})
\]
が成り立つ．
\end{lemma}
\begin{proof}
一変数論理式のコード $e$ から，
その自由変数に数詞 $\overline e$ を代入した文のコードを
返す関数 $d(e)$ を考える．
補題\ref{lem:syntax}から $d$ は原始再帰的であり，
定理\ref{thm:representability}によって
そのグラフを表す論理式 $D(x,y)$ が存在する．
\[
 \psi(x)\equiv\exists y\,
 (D(x,y)\land\theta(y))
\]
と置き，$n=\ulcorner\psi\urcorner$ とする．
$G$ を文 $\psi(\overline n)$ と定める．
そのコードは定義から $d(n)$ なので
$\ulcorner G\urcorner=d(n)$ である．
$d$ の強い表現可能性から
\[
 \mathrm{PA}\vdash
 \forall y\left(D(\overline n,y)
 \leftrightarrow y=\overline{d(n)}\right)
\]
が成立する．従って
\[
 \begin{aligned}
 \mathrm{PA}\vdash
 G&\leftrightarrow
 \exists y(D(\overline n,y)\land\theta(y))\\
 &\leftrightarrow\theta(\overline{d(n)})
 \leftrightarrow\theta(\overline{\ulcorner G\urcorner}).
 \end{aligned}
\]
これが結論である．
\end{proof}

\section{第一不完全性定理}

\begin{theorem}[Gödelの第一不完全性定理]
\label{thm:first}
$T\supseteq\mathrm{PA}$ を計算可能に公理化可能な理論とする．
$T$ が無矛盾なら，ある文 $G_T$ について
$T\not\vdash G_T$ となる．
さらに $T$ が $\omega$-無矛盾なら
$T\not\vdash\neg G_T$ も成り立つ．
\end{theorem}
\begin{proof}
対角化補題\ref{lem:diagonal}を
$\theta(x)\equiv\neg\operatorname{Prov}_T(x)$ に適用して
\[
 \mathrm{PA}\vdash
 G_T\leftrightarrow
 \neg\operatorname{Prov}_T(\overline{\ulcorner G_T\urcorner})
 \tag{1}
\]
となる $G_T$ を取る．
もし $T\vdash G_T$ なら，
補題\ref{lem:proof-numeral}から
$T\vdash\operatorname{Prov}_T(
 \overline{\ulcorner G_T\urcorner})$ となる．
一方式 (1) から
$T\vdash\neg\operatorname{Prov}_T(
 \overline{\ulcorner G_T\urcorner})$ となるので，
$T$ は矛盾する．従って $T\not\vdash G_T$ である．

さらに $\omega$-無矛盾と仮定する．
もし $T\vdash\neg G_T$ なら，
式 (1) から
$T\vdash\operatorname{Prov}_T(
\overline{\ulcorner G_T\urcorner})$ を得る．
これは
$T\vdash\exists p\,\operatorname{Proof}_T(
 p,\overline{\ulcorner G_T\urcorner})$ と同値である．
しかし既に $T\not\vdash G_T$ を示したので，
どの標準自然数 $n$ も $G_T$ の証明コードでない．
補題\ref{lem:proof-numeral}により全ての $n$ について
\[
 T\vdash\neg\operatorname{Proof}_T(
 \overline n,\overline{\ulcorner G_T\urcorner})
\]
が成り立つ．
これは $\omega$-無矛盾性の定義に反する．
従って $T\not\vdash\neg G_T$ である．
\end{proof}

\section{Rosserの改良}

\begin{theorem}[Rosserの不完全性定理]\label{thm:rosser}
$T\supseteq\mathrm{PA}$ が計算可能に公理化可能かつ
無矛盾なら，ある算術文 $R_T$ が存在して
\[
 T\not\vdash R_T,\qquad T\not\vdash\neg R_T
\]
が成り立つ．
\end{theorem}
\begin{proof}
対角化補題から
\[
 \begin{aligned}
 T\vdash R_T\leftrightarrow
 \forall p\bigl(&\operatorname{Proof}_T(
 p,\overline{\ulcorner R_T\urcorner})\\
 &\to\exists q\le p\
 \operatorname{Proof}_T(
 q,\overline{\ulcorner\neg R_T\urcorner})\bigr)
 \end{aligned}
 \tag{2}
\]
となる $R_T$ を取る．

まず $T\vdash R_T$ と仮定する．
その標準証明コードを $n$ とする．
$T$ が無矛盾なので $\neg R_T$ の標準証明は存在しない．
従って各 $q\le n$ について
$\mathrm{PA}\vdash\neg\operatorname{Proof}_T(
 \overline q,\overline{\ulcorner\neg R_T\urcorner})$
である．$q\le\overline n$ は有限個の場合に分けられるので
$\mathrm{PA}$ はその有限連言を証明できる．
一方，補題\ref{lem:proof-numeral}から
$\mathrm{PA}\vdash\operatorname{Proof}_T(
\overline n,\overline{\ulcorner R_T\urcorner})$
である．
式 (2) に $p=\overline n$ を代入すると
$T\vdash\neg R_T$ が従い，無矛盾性に反する．
従って $T\not\vdash R_T$ である．

次に $T\vdash\neg R_T$ と仮定する．
その証明コードを $n$ とする．
無矛盾性から $R_T$ の証明は一つも存在しない．
従って各 $p<n$ について
$\mathrm{PA}\vdash\neg\operatorname{Proof}_T(
 \overline p,\overline{\ulcorner R_T\urcorner})$
である．
$n$ は固定された標準自然数なので，
$\mathrm{PA}$ はこれらの有限個の数値判定を用いて
\[
 \mathrm{PA}\vdash\forall p<n\
 \neg\operatorname{Proof}_T(
 p,\overline{\ulcorner R_T\urcorner})
\]
を証明できる．
また $n$ は $\neg R_T$ の証明コードだから
\[
 \mathrm{PA}\vdash\forall p\ge n\
 \exists q\le p\
 \operatorname{Proof}_T(
 q,\overline{\ulcorner\neg R_T\urcorner})
\]
が，$q=\overline n$ を証人として成立する．
算術の線形順序によって任意の $p$ は
$p<n$ または $p\ge n$ のいずれかである．
従って $\mathrm{PA}$ は式 (2) の右辺を証明でき，
ゆえに $T\vdash R_T$ となる．
これは $T\vdash\neg R_T$ と無矛盾性に反する．
従って $T\not\vdash\neg R_T$ でもある．
\end{proof}

\section{証明可能性述語の導出可能性条件}

以下，$\ulcorner\varphi\urcorner$ を文のコードとし，
$\operatorname{Prov}_T(\ulcorner\varphi\urcorner)$ では
数詞を省略して記す．

\begin{lemma}[Hilbert--Bernays--L\"obの三条件]
\label{lem:derivability}
$T\supseteq\mathrm{PA}$ の標準的な証明可能性述語は，
任意の文 $\varphi,\psi$ について次を満たす．
\begin{enumerate}
\item $T\vdash\varphi$ なら
 $\mathrm{PA}\vdash\operatorname{Prov}_T(\ulcorner\varphi\urcorner)$．
\item
 \[
 \mathrm{PA}\vdash\operatorname{Prov}_T(
 \ulcorner\varphi\to\psi\urcorner)\to
 \bigl(\operatorname{Prov}_T(\ulcorner\varphi\urcorner)
 \to\operatorname{Prov}_T(\ulcorner\psi\urcorner)\bigr).
 \]
\item
 \[
 \mathrm{PA}\vdash\operatorname{Prov}_T(\ulcorner\varphi\urcorner)
 \to\operatorname{Prov}_T(
 \ulcorner\operatorname{Prov}_T(\ulcorner\varphi\urcorner)\urcorner).
 \]
\end{enumerate}
\end{lemma}
\begin{proof}
(1) は補題\ref{lem:proof-numeral}による．
(2) を示す．
二つの正しい証明列の末尾へmodus ponensを一回加える
原始再帰的関数 $c(p,q)$ を定義できる．
$\mathrm{PA}$ は構文検査を算術化して
\[
 \operatorname{Proof}_T(p,\ulcorner\varphi\to\psi\urcorner)
 \land\operatorname{Proof}_T(q,\ulcorner\varphi\urcorner)
 \ \to\
 \operatorname{Proof}_T(c(p,q),\ulcorner\psi\urcorner)
\]
を証明する．
証明コード $p,q$ を存在量化すれば (2) を得る．

(3) を示す．
$\operatorname{Proof}_T(p,\ulcorner\varphi\urcorner)$ の
数値検査を形式化すると，
正しい証明コード $p$ から，
その正当性を算術内で確認する証明コードを
原始再帰的に構成できる．
その末尾へ存在量化の導入を加えると，
$\operatorname{Prov}_T(\ulcorner\varphi\urcorner)$ の
$T$-証明コードを得る．
この構成を $v(p)$ とすると
$\mathrm{PA}$ は
\[
 \operatorname{Proof}_T(p,\ulcorner\varphi\urcorner)
 \to\operatorname{Proof}_T(
 v(p),
 \ulcorner\operatorname{Prov}_T(
 \ulcorner\varphi\urcorner)\urcorner)
\]
を証明できる．
$p$ を存在量化すると (3) が従う．
この証明書の符号化においても
補題\ref{lem:syntax}の構文操作と
定理\ref{thm:representability}を用いている．
\end{proof}

\section{第二不完全性定理}

\begin{definition}
$\bot$ を矛盾を表す文 $0=1$ とする．
\[
 \operatorname{Con}(T)\equiv
 \neg\operatorname{Prov}_T(\ulcorner\bot\urcorner)
\]
を理論 $T$ の無矛盾性を表す算術文という．
\end{definition}

\begin{theorem}[Gödelの第二不完全性定理]
\label{thm:second}
$T\supseteq\mathrm{PA}$ を計算可能に公理化可能で
無矛盾な理論とする．このとき
\[
 T\not\vdash\operatorname{Con}(T)
\]
が成り立つ．
\end{theorem}
\begin{proof}
第一不完全性定理の文 $G_T$ を取り，
\[
 \mathrm{PA}\vdash
 G_T\leftrightarrow
 \neg\operatorname{Prov}_T(\ulcorner G_T\urcorner)
 \tag{3}
\]
とする．
式 (3) の一方向から
$T\vdash G_T\to
 \neg\operatorname{Prov}_T(\ulcorner G_T\urcorner)$
である．
補題\ref{lem:derivability}の (1),(2) により
\[
 \mathrm{PA}\vdash
 \operatorname{Prov}_T(\ulcorner G_T\urcorner)
 \to\operatorname{Prov}_T(
 \ulcorner\neg\operatorname{Prov}_T(
 \ulcorner G_T\urcorner)\urcorner)
 \tag{4}
\]
を得る．
一方，(3) の導出可能性条件により
\[
 \mathrm{PA}\vdash
 \operatorname{Prov}_T(\ulcorner G_T\urcorner)
 \to\operatorname{Prov}_T(
 \ulcorner\operatorname{Prov}_T(
 \ulcorner G_T\urcorner)\urcorner).
 \tag{5}
\]
$T$ は命題論理の恒真式
$(A\land\neg A)\to\bot$ を証明できるので，
(1),(2) の適用によって
\[
 \mathrm{PA}\vdash
 \operatorname{Prov}_T(\ulcorner A\urcorner)
 \land\operatorname{Prov}_T(\ulcorner\neg A\urcorner)
 \to\operatorname{Prov}_T(\ulcorner\bot\urcorner)
\]
となる．
$A=\operatorname{Prov}_T(\ulcorner G_T\urcorner)$
を代入し，式 (4),(5) と合わせると
\[
 \mathrm{PA}\vdash
 \operatorname{Prov}_T(\ulcorner G_T\urcorner)
 \to\operatorname{Prov}_T(\ulcorner\bot\urcorner).
\]
従って
\[
 \mathrm{PA}\vdash
 \operatorname{Con}(T)\to
 \neg\operatorname{Prov}_T(\ulcorner G_T\urcorner).
\]
式 (3) を用いると
\[
 \mathrm{PA}\vdash\operatorname{Con}(T)\to G_T.
 \tag{6}
\]
もし $T\vdash\operatorname{Con}(T)$ なら，
$T\supseteq\mathrm{PA}$ なので式 (6) から
$T\vdash G_T$ を得る．
しかし $T$ は無矛盾であるから，
定理\ref{thm:first}によって $T\not\vdash G_T$ である．
矛盾するので
$T\not\vdash\operatorname{Con}(T)$ となる．
\end{proof}

\begin{remark}
第二不完全性定理は，証明可能性述語の
\emph{標準的な算術化}に依存する．
無矛盾性を別の非標準的な式で表現した場合，
同じ結論が無条件に成り立つとは限らない．
また定理は，より強い理論が $T$ の無矛盾性を
証明できないと主張しているわけではない．
\end{remark}

\section{到達課題の完全解答}
\noindent\textbf{Gödel文の証明不能性を実際に追う．}
\begin{proposition}[無矛盾性を使う箇所]
$T\supseteq\mathrm{PA}$ を再帰的可算公理系で
無矛盾とし，
\[
T\vdash G\leftrightarrow
\neg\operatorname{Prov}_T(\ulcorner G\urcorner)
\]
を満たす算術文 $G$ を取る．このとき
$T\not\vdash G$ である．
\end{proposition}
\begin{proof}
$T\vdash G$ と仮定する．正しい $T$-証明コードの
標準自然数 $n$ が存在するので，
証明述語の数値例の表現可能性から
$\mathrm{PA}\vdash\operatorname{Proof}_T(
\overline n,\ulcorner G\urcorner)$．
存在量化によって
$T\vdash\operatorname{Prov}_T(\ulcorner G\urcorner)$
となる．一方，上の同値式と $T\vdash G$ から
$T\vdash\neg\operatorname{Prov}_T(\ulcorner G\urcorner)$．
これは無矛盾性に反する．
\end{proof}
\noindent\textbf{到達確認．}
上の証明から，$T$ の無矛盾性だけで
$T\not\vdash\neg G$ まで導けるか．
\par\noindent\textbf{解答．}
このGödel文に対しては一般に導けない．
否定の証明不能性には $\omega$-無矛盾性など
追加の仮定が必要である．
単なる無矛盾性から独立な文を得るには
Rosserの改良を用いる．
\end{document}